\chapter*{Introduction générale}
\addcontentsline{toc}{chapter}{Introduction générale}

Dans le contexte actuel de la transformation digitale et de la compétitivité accrue sur le marché du travail, la gestion des ressources humaines est devenue un levier stratégique majeur pour la performance des organisations. Le capital humain, considéré comme l'atout le plus précieux d'une entreprise, nécessite une attention particulière, notamment en ce qui concerne la rétention des talents et l'optimisation de la satisfaction au travail. C'est dans cette optique que s'inscrit le domaine du \textit{HR Analytics}, une discipline qui allie la science des données et la gestion des ressources humaines pour transformer des volumes massifs d'informations en insights exploitables.

Le présent projet, intitulé \textbf{« HR Analytics - Employee Attrition \& Performance »}, a été réalisé dans le cadre du module de Business Intelligence. Son objectif principal est de concevoir et de mettre en œuvre une solution décisionnelle permettant d'analyser les facteurs déterminants de l'attrition des employés et d'évaluer les niveaux de performance au sein d'une organisation. En s'appuyant sur des données empiriques, cette étude vise à identifier les motifs récurrents de départ et à proposer des indicateurs de performance clés (KPIs) pour guider les prises de décisions managériales.

La démarche adoptée suit le cycle de vie classique d'un projet de Business Intelligence, allant de l'expression des besoins à la visualisation finale des résultats. Ce rapport est ainsi structuré en six chapitres principaux :

Le premier chapitre pose le cadre général de l'étude en présentant le contexte du projet, la problématique adressée ainsi que les objectifs visés.

Le deuxième chapitre est dédié à l'analyse fonctionnelle des besoins et à l'exploration préliminaire des données, étape cruciale pour comprendre la structure et la qualité des informations disponibles.

Le troisième chapitre expose la phase de modélisation dimensionnelle, où nous détaillerons la conception de l'entrepôt de données, notamment le schéma en étoile retenu pour cette analyse.

Le quatrième chapitre décrit le processus ETL (Extract, Transform, Load), explicitant les différentes transformations appliquées aux données pour garantir leur intégrité et leur pertinence.

Le cinquième chapitre présente les résultats de l'analyse à travers des tableaux de bord interactifs et propose une interprétation approfondie des indicateurs obtenus.

Enfin, une conclusion générale synthétise les apports du projet et ouvre des perspectives sur des évolutions futures, notamment l'intégration de modèles prédictifs.
